Start with the squared-up dynamic multiplier
$\bmat{v_1(t) \\ v_2(t)}=\Psi\bmat{z_\Delta(t) \\ w_\Delta(t)}$,
\begin{equation}\label{eqn:multiplier_untransformed}
    \bmat{\mcl{T}_\Psi\dot{\mbf{x}}_\Psi(t)\\v_1(t)\\v_2(t)}
    =
    \bmat{
        \mcl{A}_\Psi & \mcl{B}_\Psi^z & \mcl{B}_\Psi^w\\
        \mcl{C}_\Psi^z & \mcl{D}_\Psi^{z} & \mcl{D}_\Psi^{zw} \\
        \mcl{C}_\Psi^w & \mcl{D}_\Psi^{wz} & \mcl{D}_\Psi^w
    }
    \bmat{\mbf{x}_\Psi(t) \\ z_\Delta(t) \\ w_\Delta(t)}.
\end{equation}
To square up the interconnection, introduce a full-row-rank map $\mcl{N}$ from
the squared coordinates to the original coordinates, and choose any right
inverse $\mcl{E}$ satisfying
\begin{equation}\label{eqn:squareup_maps}
    \mcl{N}\mcl{E}=I.
\end{equation}
The choice $\mcl{E}=\mcl{N}^*$ is the usual zero-padding choice when $\mcl{N}$
is an orthogonal selector, but nothing below requires this. The original plant
sees the projected input $\mcl{N}w_\Delta$, while its original uncertainty
output is embedded into the squared coordinates by $\mcl{E}$. The signal route is
\begin{equation}\label{eqn:squareup_signal_route}
    w_{\Delta,\square}
    \xrightarrow{\mcl{N}}
    w_{\Delta,\mathrm{o}}
    \xrightarrow{G}
    z_{\Delta,\mathrm{o}}
    \xrightarrow{\mcl{E}}
    z_{\Delta,\square}
    \xrightarrow{\Psi^\dagger_\square}
    w_{\Delta,\square} .
\end{equation}
The maps $\mcl{N}$ and $\mcl{E}$ appear in the plant blocks. If these maps are
dynamic, their states are appended to the augmented realization; the algebraic
loop below is then obtained from their direct feedthrough operators. Thus the
squared-up generalized plant is
\begin{equation}\label{eqn:squaredup_gp}
    \bmat{\mcl{T}\dot{\mbf{x}}\\z_{\Delta,\square}\\z}
    =
    \bmat{
        \mcl{A} & \mcl{B}_\Delta\mcl{N} & \mcl{B}_w\\
        \mcl{E}\mcl{C}_\Delta
        & \mcl{E}\mcl{D}_{\Delta}\mcl{N}
        & \mcl{E}\mcl{D}_{\Delta w}\\
        \mcl{C}_z & \mcl{D}_{z\Delta}\mcl{N} & \mcl{D}_{zw}
    }
    \bmat{\mbf{x}\\w_{\Delta,\square}\\w}.
\end{equation}

Use the abbreviations
\begin{equation}
    \begin{gathered}
    \mcl{C}_\Delta^\square:=\mcl{E}\mcl{C}_\Delta,\qquad
    \mcl{D}_\Delta^\square:=\mcl{E}\mcl{D}_{\Delta}\mcl{N},\qquad
    \mcl{D}_{\Delta w}^\square:=\mcl{E}\mcl{D}_{\Delta w},\\
    \mcl{B}_\Delta^\square:=\mcl{B}_\Delta\mcl{N},\qquad
    \mcl{D}_{z\Delta}^\square:=\mcl{D}_{z\Delta}\mcl{N},\qquad
    \mbf{x}^a:=\bmat{\mbf{x}\\\mbf{x}_\Psi},\qquad
    \mcl{T}^a:=\bmat{\mcl{T}&0\\0&\mcl{T}_\Psi}.
    \end{gathered}
\end{equation}
Substituting the squared-up plant signals into the multiplier realization and
stacking the plant and multiplier states gives the open augmented realization
with state $\mbf{x}^a=\bmat{\mbf{x}\\\mbf{x}_\Psi}$. Before closing the
algebraic loop, this gives
\begin{equation}\label{eqn:aug_untransformed}
    \bmat{
        \mcl{T}^a\dot{\mbf{x}}^a(t)\\
        v_1(t)\\
        v_2(t)\\
        z(t)
    }
    =
    \bmat{
        \mcl{A}^a & \mcl{B}^a_{\Delta} & \mcl{B}^a_{w}\\
        \mcl{C}^a_1 & \mcl{D}_{11}^a & \mcl{D}_{12}^a\\
        \mcl{C}^a_2 & \mcl{D}_{21}^a & \mcl{D}_{22}^a\\
        \mcl{C}_z^a & \mcl{D}_{z\Delta}^\square & \mcl{D}_{zw}
    }
    \bmat{\mbf{x}^a(t)\\w_{\Delta,\square}(t)\\w(t)},
\end{equation}
where
\begin{align}
    \mcl{A}^a
    &:=
    \bmat{
        \mcl{A} & 0\\
        \mcl{B}_\Psi^z\mcl{C}_\Delta^\square & \mcl{A}_\Psi
    },
    &
    \mcl{B}^a_{\Delta}
    &:=
    \bmat{
        \mcl{B}_\Delta^\square\\
        \mcl{B}_\Psi^z\mcl{D}_\Delta^\square+\mcl{B}_\Psi^w
    },
    &
    \mcl{B}^a_{w}
    &:=
    \bmat{
        \mcl{B}_w\\
        \mcl{B}_\Psi^z\mcl{D}_{\Delta w}^\square
    },
    \nonumber\\
    \mcl{C}^a_1
    &:=
    \bmat{\mcl{D}_\Psi^z\mcl{C}_\Delta^\square & \mcl{C}_\Psi^z},
    &
    \mcl{D}^a_{11}
    &:=
    \mcl{D}_\Psi^z\mcl{D}_\Delta^\square+\mcl{D}_\Psi^{zw},
    &
    \mcl{D}_{12}^a
    &:=
    \mcl{D}_\Psi^z\mcl{D}_{\Delta w}^\square,
    \nonumber\\
    \mcl{C}^a_2
    &:=
    \bmat{\mcl{D}_\Psi^{wz}\mcl{C}_\Delta^\square & \mcl{C}_\Psi^w},
    &
    \mcl{D}_{21}^a
    &:=
    \mcl{D}_\Psi^{wz}\mcl{D}_\Delta^\square+\mcl{D}_\Psi^w,
    &
    \mcl{D}_{22}^a
    &:=
    \mcl{D}_\Psi^{wz}\mcl{D}_{\Delta w}^\square,
    \nonumber\\
        \mcl{C}_z^a
    &:= 
    \bmat{\mcl{C}_z & 0}.
\end{align}
The algebraic loop is entirely contained in the $v_2$-row:
\begin{equation}\label{eqn:v2_loop_row}
    v_2
    =
    \mcl{C}^a_2\mbf{x}^a
    +\mcl{D}^a_{21}w_{\Delta,\square}
    +\mcl{D}_{22}^a w.
\end{equation}

Assume that $\mcl{D}_{21}^a$ is invertible. Then
\begin{equation}\label{eqn:wdelta_from_v2}
    w_{\Delta,\square}
    =
    -(\mcl{D}_{21}^a)^{-1}\mcl{C}^a_2\mbf{x}^a
    +(\mcl{D}_{21}^a)^{-1}v_2
    -(\mcl{D}_{21}^a)^{-1}\mcl{D}_{22}^a w.
\end{equation}
Hence the change of variables
\begin{equation}\label{eqn:loop_congruence}
    \bmat{\mbf{x}^a\\w_{\Delta,\square}\\w}
    =
    \underbrace{
    \bmat{
        I & 0 & 0\\
        -(\mcl{D}_{21}^a)^{-1}\mcl{C}^a_2
        & (\mcl{D}_{21}^a)^{-1}
        & -(\mcl{D}_{21}^a)^{-1}\mcl{D}_{22}^a\\
        0 & 0 & I
    }}_{\mcl{X}_\Delta}
    \bmat{\mbf{x}^a\\v_2\\w}
\end{equation}
closes the loop. This is the congruence transformation: any quadratic form
$L$ in the variables $\bmat{\mbf{x}^a\\w_{\Delta,\square}\\w}$ becomes
\begin{equation}\label{eqn:lmi_congruence}
    L
    \quad\longmapsto\quad
    \mcl{X}_\Delta^*L\mcl{X}_\Delta
\end{equation}
in the variables $\bmat{\mbf{x}^a\\v_2\\w}$.
The inverse change of variables is just the original loop row
\eqref{eqn:v2_loop_row}:
\begin{equation}\label{eqn:loop_congruence_inverse}
    \bmat{\mbf{x}^a\\v_2\\w}
    =
    \underbrace{
    \bmat{
        I & 0 & 0\\
        \mcl{C}^a_2 & \mcl{D}_{21}^a & \mcl{D}_{22}^a\\
        0 & 0 & I
    }}_{\mcl{X}_\Delta^{-1}}
    \bmat{\mbf{x}^a\\w_{\Delta,\square}\\w}.
\end{equation}
% We now relate this squared-up condition to the original, non-squared-up
% interconnection.
% The order of $\mcl{N}$ and $\mcl{F}$ is determined by the signal space in
% which the loop is written. Let $\mcl{W}_{\Delta,\mathrm{o}}$ denote the
% original uncertainty-input space and let
% $\mcl{W}_{\Delta,\square}$ denote the squared-up uncertainty-input space. The
% operator $\mcl{N}:\mcl{W}_{\Delta,\square}\to
% \mcl{W}_{\Delta,\mathrm{o}}$ projects squared-up inputs back to the original
% plant input. The operator
% \begin{equation}
%     \mcl{F}
%     =
%     (\mcl{D}_\Psi^w)^{-1}
%     \mcl{D}_\Psi^{wz}
%     \mcl{E}\mcl{D}_\Delta
% \end{equation}
% goes in the opposite direction:
% \begin{equation}
%     \mcl{F}:\mcl{W}_{\Delta,\mathrm{o}}
%     \longrightarrow
%     \mcl{W}_{\Delta,\square}.
% \end{equation}
% Indeed, starting from an original uncertainty input
% $w_{\Delta,\mathrm{o}}$, the direct feedthrough route is
% \begin{equation}
%     w_{\Delta,\mathrm{o}}
%     \xrightarrow{\mcl{D}_\Delta}
%     z_{\Delta,\mathrm{o}}
%     \xrightarrow{\mcl{E}}
%     z_{\Delta,\square}
%     \xrightarrow{\mcl{D}_\Psi^{wz}}
%     v_2
%     \xrightarrow{(\mcl{D}_\Psi^w)^{-1}}
%     w_{\Delta,\square}.
% \end{equation}
% Thus, when the loop is written in the original coordinate
% $w_{\Delta,\mathrm{o}}$, one loop pass is
% $w_{\Delta,\mathrm{o}}\mapsto\mcl{N}\mcl{F}w_{\Delta,\mathrm{o}}$, giving
% $I+\mcl{N}\mcl{F}$. When the loop is written in the squared coordinate
% $w_{\Delta,\square}$, the plant first sees
% $\mcl{N}w_{\Delta,\square}$, and one loop pass is
% $w_{\Delta,\square}\mapsto\mcl{F}\mcl{N}w_{\Delta,\square}$, giving
% $I+\mcl{F}\mcl{N}$.

% \begin{theorem}[Well-posedness under square-up]\label{thm:squareup_wellposed}
% Assume $\mcl{D}_\Psi^w$ is invertible. If the original non-squared-up loop
% operator
% \begin{equation}\label{eqn:nonpadded_loop_untransformed}
%     \mcl{L}_{\mathrm{ns}}
%     :=
%     I+\mcl{N}\mcl{F}
%     =
%     I
%     +
%     \mcl{N}
%     (\mcl{D}_\Psi^w)^{-1}
%     \mcl{D}_\Psi^{wz}
%     \mcl{E}
%     \mcl{D}_\Delta
% \end{equation}
% is invertible, then the squared-up loop operator $I+\mcl{F}\mcl{N}$ is
% invertible. Consequently the squared-up algebraic loop is well posed, since
% \begin{equation}\label{eqn:squared_wellposed_factor}
%     \mcl{D}_{21}^a
%     =
%     \mcl{D}_\Psi^w
%     \left(I+\mcl{F}\mcl{N}\right).
% \end{equation}
% \end{theorem}

% \begin{proof}
% Using $\mcl{D}_\Delta^\square=\mcl{E}\mcl{D}_\Delta\mcl{N}$ from
% \eqref{eqn:squareup_maps}, the loop feedthrough in the augmented squared-up
% system factors as
% \[
%     \mcl{D}_{21}^a
%     =
%     \mcl{D}_\Psi^{wz}\mcl{D}_\Delta^\square+\mcl{D}_\Psi^w
%     =
%     \mcl{D}_\Psi^w
%     \left(I+\mcl{F}\mcl{N}\right).
% \]
% If $\mcl{L}_{\mathrm{ns}}=I+\mcl{N}\mcl{F}$ is invertible, then
% \begin{equation}\label{eqn:squared_inverse_from_original}
%     \left(I+\mcl{F}\mcl{N}\right)^{-1}
%     =
%     I-\mcl{F}
%     \left(I+\mcl{N}\mcl{F}\right)^{-1}
%     \mcl{N}.
% \end{equation}
% Thus $I+\mcl{F}\mcl{N}$ is invertible, and so is $\mcl{D}_{21}^a$ by
% \eqref{eqn:squared_wellposed_factor}. Equivalently, if
% $w_\square\in\ker(I+\mcl{F}\mcl{N})$, then applying $\mcl{N}$ gives
% $\mcl{L}_{\mathrm{ns}}\mcl{N}w_\square=0$, so $\mcl{N}w_\square=0$; the
% original equation then gives $w_\square=0$. Thus the added completion
% directions do not introduce new algebraic-loop singularities.
% \end{proof}

\begin{example}[Zero-padding completion]
For the trivial zero-padding completion the well-posedness implication is visible block by
block. Let
\begin{equation}
    \mcl{N}=\bmat{I&0},
    \qquad
    \mcl{E}=\mcl{N}^*=\bmat{I\\0},
    \qquad
    w_{\Delta,\square}=\bmat{w_{\Delta,\mathrm{o}}\\w_{\Delta,\mathrm{p}}},
\end{equation}
so that $\mcl{N}w_{\Delta,\square}=w_{\Delta,\mathrm{o}}$, returns the original uncertainty input. Suppose $(I +\mcl{L})$ denotes the 
well-posedness condition. $\mcl{L}$ can be separated in a part $\mcl{F}$ and $\mcl{N}$.
Since $\mcl{F}$ can be separated in parts acting on the original coordinates
and the padded coordinates, write
\begin{equation}
    \mcl{F}=\bmat{\mcl{F}_{\mathrm{o}}\\\mcl{F}_{\mathrm{p}}}.
\end{equation}
Then the squared-up, and non-squared-up algebraic loop operators are
\begin{equation}
    I+\mcl{F}\mcl{N}
    =
    \bmat{
        I+\mcl{F}_{\mathrm{o}} & 0\\
        \mcl{F}_{\mathrm{p}} & I
    },
    \qquad
    I+\mcl{N}\mcl{F}
    =
    I+\mcl{F}_{\mathrm{o}},
\end{equation}
respectively.
Thus the squared-up loop is block lower triangular: the padded directions add
only an identity block, and the only possible singularity is the original
non-squared-up loop $I+\mcl{F}_{\mathrm{o}}$.
\end{example}



% The dual-side construction should be read in the same way as
% Theorem~\ref{thm:dualscaledIO}: we do not first postulate a separate dual
% multiplier. Instead, take the adjoint of the open augmentation
% \eqref{eqn:aug_untransformed} and then apply the corresponding dual
% congruence.

% The adjoint of the squared-up generalized plant \eqref{eqn:squaredup_gp} is
% \begin{equation}\label{eqn:dual_squaredup_gp}
%     \bmat{
%         \mcl{T}^*\dot{\underline{\mbf{x}}}(t)\\
%         \underline{w}_{\Delta,\square}^{G}(t)\\
%         \underline{w}
%     }
%     =
%     \bmat{
%         \mcl{A}^*
%         & (\mcl{C}_\Delta^\square)^*
%         & \mcl{C}_z^*\\
%         (\mcl{B}_\Delta^\square)^*
%         & (\mcl{D}_\Delta^\square)^*
%         & (\mcl{D}_{z\Delta}^\square)^*\\
%         \mcl{B}_w^*
%         & (\mcl{D}_{\Delta w}^\square)^*
%         & \mcl{D}_{zw}^*
%     }
%     \bmat{
%         \underline{\mbf{x}}(t)\\
%         \underline{z}_{\Delta,\square}(t)\\
%         \underline{z}(t)
%     }.
% \end{equation}

% The dual dynamic multiplier is the adjoint of
% \eqref{eqn:multiplier_untransformed}:
% \begin{equation}\label{eqn:dual_dynamic_multiplier}
%     \bmat{
%         \mcl{T}_\Psi^*\dot{\underline{\mbf{x}}}_\Psi(t)\\
%         \underline{z}_{\Delta,\square}(t)\\
%         \underline{w}_{\Delta,\square}^{\Psi}(t)
%     }
%     =
%     \bmat{
%         \mcl{A}_\Psi^*
%         & (\mcl{C}_\Psi^z)^*
%         & (\mcl{C}_\Psi^w)^*\\
%         (\mcl{B}_\Psi^z)^*
%         & (\mcl{D}_\Psi^z)^*
%         & (\mcl{D}_\Psi^{wz})^*\\
%         (\mcl{B}_\Psi^w)^*
%         & (\mcl{D}_\Psi^{zw})^*
%         & (\mcl{D}_\Psi^w)^*
%     }
%     \bmat{
%         \underline{\mbf{x}}_\Psi(t)\\
%         \underline{v}_1(t)\\
%         \underline{v}_2(t)
%     }.
% \end{equation}

% The two dual systems are augmented by feeding the multiplier output
% $\underline{z}_{\Delta,\square}$ into the adjoint generalized plant. From
% \eqref{eqn:dual_dynamic_multiplier},
% \begin{equation}
%     \underline{z}_{\Delta,\square}
%     =
%     (\mcl{B}_\Psi^z)^*\underline{\mbf{x}}_\Psi
%     +(\mcl{D}_\Psi^z)^*\underline{v}_1
%     +(\mcl{D}_\Psi^{wz})^*\underline{v}_2 .
% \end{equation}
% Substituting this into the state equation of
% \eqref{eqn:dual_squaredup_gp} and keeping the multiplier state equation gives
% \begin{equation}
%     \begin{aligned}
%     \bmat{
%         \mcl{T}^*\dot{\underline{\mbf{x}}}\\
%         \mcl{T}_\Psi^*\dot{\underline{\mbf{x}}}_\Psi
%     }
%     &=
%     \bmat{
%         \mcl{A}^*
%         & (\mcl{C}_\Delta^\square)^*(\mcl{B}_\Psi^z)^*\\
%         0
%         & \mcl{A}_\Psi^*
%     }
%     \bmat{\underline{\mbf{x}}\\\underline{\mbf{x}}_\Psi}
%     +
%     \bmat{
%         (\mcl{C}_\Delta^\square)^*(\mcl{D}_\Psi^z)^*\\
%         (\mcl{C}_\Psi^z)^*
%     }\underline{v}_1\\
%     &\quad+
%     \bmat{
%         (\mcl{C}_\Delta^\square)^*(\mcl{D}_\Psi^{wz})^*\\
%         (\mcl{C}_\Psi^w)^*
%     }\underline{v}_2
%     +
%     \bmat{\mcl{C}_z^*\\0}\underline{z}.
%     \end{aligned}
% \end{equation}
% The dual uncertainty input is the sum of the plant and multiplier
% contributions:
% \begin{equation}
%     \begin{aligned}
%     \underline{w}_{\Delta,\square}
%     &:=
%     \underline{w}_{\Delta,\square}^{G}
%     +
%     \underline{w}_{\Delta,\square}^{\Psi}\\
%     &=
%     \left(
%     (\mcl{B}_\Delta^\square)^*
%     \quad
%     (\mcl{D}_\Delta^\square)^*(\mcl{B}_\Psi^z)^*
%     +(\mcl{B}_\Psi^w)^*
%     \right)
%     \bmat{\underline{\mbf{x}}\\\underline{\mbf{x}}_\Psi}\\
%     &\quad+
%     \left((\mcl{D}_\Delta^\square)^*(\mcl{D}_\Psi^z)^*
%     +(\mcl{D}_\Psi^{zw})^*\right)\underline{v}_1
%     +
%     \left((\mcl{D}_\Delta^\square)^*(\mcl{D}_\Psi^{wz})^*
%     +(\mcl{D}_\Psi^w)^*\right)\underline{v}_2
%     +(\mcl{D}_{z\Delta}^\square)^*\underline{z}.
%     \end{aligned}
% \end{equation}
% Likewise, after the same substitution, the adjoint external channel is
% \begin{equation}
%     \begin{aligned}
%     \underline{w}
%     &=
%     \left(
%     \mcl{B}_w^*
%     \quad
%     (\mcl{D}_{\Delta w}^\square)^*(\mcl{B}_\Psi^z)^*
%     \right)
%     \bmat{\underline{\mbf{x}}\\\underline{\mbf{x}}_\Psi}\\
%     &\quad+
%     (\mcl{D}_{\Delta w}^\square)^*(\mcl{D}_\Psi^z)^*\underline{v}_1
%     +(\mcl{D}_{\Delta w}^\square)^*(\mcl{D}_\Psi^{wz})^*\underline{v}_2
%     +\mcl{D}_{zw}^*\underline{z}.
%     \end{aligned}
% \end{equation}
% With the augmented dual state
% $\underline{\mbf{x}}^a:=\bmat{\underline{\mbf{x}}\\
% \underline{\mbf{x}}_\Psi}$, the resulting realization is
% \begin{equation}\label{eqn:dual_augmented_unclosed}
%     \bmat{
%         (\mcl{T}^a)^*\dot{\underline{\mbf{x}}}^a(t)\\
%         \underline{w}_{\Delta,\square}(t)\\
%         \underline{w}(t)
%     }
%     =
%     \bmat{
%         (\mcl{A}^a)^*
%         & (\mcl{C}_1^a)^*
%         & (\mcl{C}_2^a)^*
%         & (\mcl{C}_z^a)^*\\
%         (\mcl{B}_\Delta^a)^*
%         & (\mcl{D}_{11}^a)^*
%         & (\mcl{D}_{21}^a)^*
%         & (\mcl{D}_{z\Delta}^\square)^*\\
%         (\mcl{B}_w^a)^*
%         & (\mcl{D}_{12}^a)^*
%         & (\mcl{D}_{22}^a)^*
%         & \mcl{D}_{zw}^*
%     }
%     \bmat{
%         \underline{\mbf{x}}^a(t)\\
%         \underline{v}_1(t)\\
%         \underline{v}_2(t)\\
%         \underline{z}(t)
%     }.
% \end{equation}
% Thus \eqref{eqn:dual_augmented_unclosed} is exactly the adjoint of the open
% primal augmentation \eqref{eqn:aug_untransformed}; the later dual congruence
% closes the $\underline{v}_2$ loop.

\begin{definition}
    Let $\mcl{V}=\mcl{V}^*\in\Pi_4$ be partitioned as
    \begin{equation*}
        \mcl{V}
        =
        \bmat{\mcl{Q} & \mcl{S}\\\mcl{S}^* & \mcl{R}}.
    \end{equation*}
    Then $\mcl{V}$ is called a strict PN-PI operator if there exists
    $\epsilon>0$ such that
    \begin{equation*}
        \mcl{Q}\succeq\epsilon I,
        \qquad
        \mcl{R}\preceq-\epsilon I.
    \end{equation*}
\end{definition}

\begin{lemma}\label{lem:primdualPNPI}
    Let $\mcl{V}=\mcl{V}^*\in\Pi_4$ and
    $\bar{\mcl{V}}=\bar{\mcl{V}}^*\in\Pi_4$ denote the static primal and dual
    IQC multipliers. Then $\mcl{V}$ is a strict PN-PI operator if and only if
    $\bar{\mcl{V}}$ is a strict PN-PI operator.
\end{lemma}
\begin{proof}
    See Appendix~\ref{proof:primdualPNPI}.
\end{proof}

\begin{example}[Dual multiplier for the zero-padded LTI-parametric IQC]
    In this example we assume that the scaled duality theorem holds for
    dynamic scalings with invertible feedthrough, and we use the square-up
    construction to apply it to the non-invertible factor.  
    Consider the case of LTI parametric uncertainty. The primal IQC multiplier
    is given by
    \begin{equation}
        \int_0^T
        \ip{\bmat{z_\Delta(t) \\ w_\Delta(t)}}
        {
        \bmat{\Psi_v^* & 0 \\ 0 &\Psi_v^*}
        \bmat{\mcl{Q} & \mcl{S} \\ \mcl{S}^* & \mcl{R}}
        \bmat{\Psi_v & 0 \\ 0 &\Psi_v}
        \bmat{z_\Delta(t) \\ w_\Delta(t)}
        }dt
        \geq 0,
    \end{equation}
    where $\Psi_v$ is stable, $\mcl{Q}\succ0$, $\mcl{R}\prec0$, and
    $\mcl{S}=-\mcl{S}^*$. Usually $\Psi_v$ is tall, with
    $\Psi_v=\bmat{I\\\Psi_N}$, where $N$ denotes the spectral order of the
    dynamic multiplier. Complete this factor by
    \begin{equation}
        \Psi_{v,\square}
        =
        \bmat{I & -\Psi_N^*\\ \Psi_N & I},
        \qquad
        \Psi_\square
        =
        \bmat{\Psi_{v,\square} & 0 \\ 0 & \Psi_{v,\square}}.
    \end{equation}
    The squared up primal IQC then becomes
    \begin{equation}
        \int_0^T
        \ip{\bmat{z_{\Delta,\square}(t) \\ w_{\Delta,\square}(t)}}
        {
        \bmat{\Psi_{v,\square}^* & 0 \\ 0 &\Psi_{v,\square}^*}
        \bmat{\mcl{Q} & \mcl{S} \\ \mcl{S}^* & \mcl{R}}
        \bmat{\Psi_{v,\square} & 0 \\ 0 &\Psi_{v,\square}}
        \bmat{z_{\Delta,\square}(t) \\ w_{\Delta,\square}(t)}
        }dt
        \geq 0,
    \end{equation}
    with $z_{\Delta,\square} = \bmat{z_\Delta^\top & 0^\top}^\top$ and $w_{\Delta,\square} = \bmat{w_\Delta^\top & 0^\top}^\top$ denoting the zero-padded signals. The multiplication operator $\mcl{M}$ admits a $J$-factorization
    \begin{equation}
        \mcl{M}
        =
        \bmat{\mcl{Q} & \mcl{S} \\ \mcl{S}^* & \mcl{R}} 
        =
        \mcl{V}^*J\mcl{V},
        \qquad
        \mcl{V}
        :=
        \bmat{\mcl{V}_{11} & \mcl{V}_{12}\\0 & \mcl{V}_{22}}, \qquad
        J:=
        \bmat{I & 0\\0 & -I},
    \end{equation}
    where $\mcl{S}_{11}$ and $\mcl{S}_{22}$ are invertible. Hence the squared
    IQC multiplier is
    \begin{equation}
        \Pi_\square
        =
        \Psi_\square^*\mcl{M}\Psi_\square
        =
        \left(\mcl{V}\Psi_\square\right)^*
        J
        \left(\mcl{V}\Psi_\square\right).
    \end{equation}
    The corresponding square primal scaling is the $J$-factor applied after
    the square dynamic filter:
    \begin{equation}
        \mcl{L}
        :=
        \mcl{V}\Psi_\square
        =
        \bmat{
            \mcl{S}_{11}\Psi_{v,\square}
            &
            \mcl{S}_{12}\Psi_{v,\square}\\
            0
            &
            \mcl{S}_{22}\Psi_{v,\square}
        }.
    \end{equation}
    Theorem~\ref{thm:dualscaledIO} uses the inverse signal route. Thus define
    \begin{equation}
        \bmat{z_\Delta\\w_\Delta}
        =
        \mcl{L}^{-1}
        \bmat{\tilde z_\Delta\\\tilde w_\Delta}.
    \end{equation}
    If
    \begin{equation}
        \mcl{A}_L:=\mcl{S}_{11}\Psi_{v,\square},\qquad
        \mcl{B}_L:=\mcl{S}_{12}\Psi_{v,\square},\qquad
        \mcl{D}_L:=\mcl{S}_{22}\Psi_{v,\square},
    \end{equation}
    then
    \begin{equation}
        \mcl{L}^{-1}
        =
        \bmat{
            \mcl{A}_L^{-1}
            &
            -\mcl{A}_L^{-1}\mcl{B}_L\mcl{D}_L^{-1}\\
            0
            &
            \mcl{D}_L^{-1}
        }.
    \end{equation}
    Applying \eqref{eqn:dualscaling} gives the dual I/O scaling
    \begin{equation}
        \underline{\mcl{L}}^{-1}
        =
        \bmat{
            \mcl{D}_L^{-*}
            &
            \mcl{D}_L^{-*}\mcl{B}_L^*\mcl{A}_L^{-*}\\
            0
            &
            \mcl{A}_L^{-*}
        }^{-1}.
    \end{equation}
    Therefore
    \begin{equation}
        \underline{\mcl{L}}
        =
        \underbrace{\bmat{
            \mcl{S}_{22}^{-*}
            &
            \mcl{S}_{22}^{-*}\mcl{S}_{12}^*\mcl{S}_{11}^{-*}\\
            0
            &
            \mcl{S}_{11}^{-*}
        }}_{\underline{\mcl{V}}}
        \underbrace{\bmat{
            \Psi_{v,\square}^{-*} & 0\\
            0 & \Psi_{v,\square}^{-*}
        }}_{\underline{\Psi}_\square}.
    \end{equation}
    The corresponding square dual multiplier is
    \begin{equation}
        \underline{\Pi}_\square
        =
        \underline{\mcl{L}}^*J\underline{\mcl{L}}
        =
        \underline{\Psi}_\square^*
        \underbrace{\underline{\mcl{V}}^*J\underline{\mcl{V}}}
        _{\underline{\mcl{M}}}
        \underline{\Psi}_\square .
    \end{equation}
    Equivalently, use the same static dual construction as in
    Lemma~\ref{lem:primdualPNPI}. If
    \begin{equation}
        \mcl{M}
        =
        \bmat{
            \mcl{Q} & \mcl{S}\\
            \mcl{S}^* & \mcl{R}
        },
        \qquad
        \mcl{K}
        :=
        \bmat{
            -\mcl{R} & \mcl{S}^*\\
            \mcl{S} & -\mcl{Q}
        },
    \end{equation}
    then the square static dual multiplier is
    \begin{equation}
        \underline{\mcl{M}}
        :=
        \mcl{K}^{-1}
        =
        \bmat{
            \underline{\mcl{Q}}
            &
            \underline{\mcl{S}}\\
            \underline{\mcl{S}}^*
            &
            \underline{\mcl{R}}
        }.
    \end{equation}
    Applying the block inverse formula with the Schur complement
    \begin{equation}
        \Theta
        :=
        \mcl{S}^*\mcl{Q}^{-1}\mcl{S}
        -\mcl{R}
    \end{equation}
    gives
    \begin{equation}
        \underline{\mcl{Q}}=\Theta^{-1} \succ 0 ,
        \qquad
        \underline{\mcl{S}}
        =
        \Theta^{-1}\mcl{S}^*
        \mcl{Q}^{-1},
    \end{equation}
    \begin{equation}
        \underline{\mcl{R}}
        =
        -\left(
            \mcl{Q}
            +
            \mcl{S}(-\mcl{R})^{-1}
            \mcl{S}^*
        \right)^{-1}\prec 0.
    \end{equation}
    Thus the square dual IQC condition is
    \begin{equation}
        \int_0^T
        \ip{\bmat{\underline{z}_\Delta(t) \\ \underline{w}_\Delta(t)}}
        {
        \underline{\Pi}_\square
        \bmat{\underline{z}_\Delta(t) \\ \underline{w}_\Delta(t)}
        }dt
        \leq 0.
    \end{equation}
    To unsquare this condition, let
    \begin{equation}
        \mcl{N}=\bmat{I&0},
        \qquad
        \mcl{N}^*=\bmat{I\\0}.
    \end{equation}
    The adjoint square-up embeds the non-squared dual uncertainty signals as
    \begin{equation}
        \bmat{\underline{z}_{\Delta,\square}\\\underline{w}_{\Delta,\square}}
        =
        \underbrace{\bmat{\mcl{N}^*&0\\0&\mcl{N}^*}}
        _{\underline{\mcl{E}}}
        \bmat{\underline{z}_{\Delta,\mathrm{o}}\\\underline{w}_{\Delta,\mathrm{o}}}.
    \end{equation}
    Hence the non-squared dual multiplier is
    \begin{equation}
        \underline{\Pi}_{\mathrm{o}}
        =
        \underline{\mcl{E}}^*
        \underline{\Pi}_\square
        \underline{\mcl{E}}
        =
        \left(\underline{\mcl{L}}\underline{\mcl{E}}\right)^*
        J
        \left(\underline{\mcl{L}}\underline{\mcl{E}}\right).
    \end{equation}
    For the zero-padding completion used here,
    \begin{equation}
        \underline{\Psi}_v=\Psi_{v,\square}^{-*}\mcl{N}^*
        =
        \bmat{(I +\Psi_N^*\Psi_N)^{-1} & -(I +\Psi_N^*\Psi_N)^{-1}\Psi_N^{*} \\
         \Psi_N(I+\Psi_N^*\Psi_N)^{-1} & I - \Psi_N(I+\Psi_N^*\Psi_N)^{-1}\Psi_N^*}\bmat{I\\0}
        =
        (I+\Psi_N^*\Psi_N)^{-1}\Psi_v.
    \end{equation}
    
    The unsquared dual IQC is then
    \begin{equation}
        \int_0^T
        \ip{\bmat{\underline{z}_{\Delta}(t) \\
        \underline{w}_{\Delta}(t)}}
        {
        \bmat{\underline{\Psi}_v^* & 0 \\ 0 &\underline{\Psi}_v^*}\bmat{
            \underline{\mcl{Q}}
            &
            \underline{\mcl{S}}\\
            \underline{\mcl{S}}^*
            &
            \underline{\mcl{R}}
        }\bmat{\underline{\Psi}_v & 0 \\ 0 & \underline{\Psi}_v}
        \bmat{\underline{z}_{\Delta}(t) \\
        \underline{w}_{\Delta}(t)}
        }dt
        \leq 0.
    \end{equation}
    which is also in the class of LTI parametric IQCs. Since all operations are invertible, we can start with
    the dual system and end up with a multiplier that satisfies the primal IQC. A convenient
    choice of the dynamic multiplier is all--pass basis function (12B) in \cite{veenman2016robust}
    \begin{equation}
    \Psi_v
    =
    \bmat{I & \frac{s + \rho}{s-\rho} & \frac{(s + \rho)^2}{(s-\rho)^2} & \cdots & \frac{(s + \rho)^v}{(s-\rho)^v}}^\top,
    \qquad
    \end{equation}
    with $\rho < 0$. This choice is convenient as $\Psi_N^*\Psi_N$ is a scalar. Indeed,
    \begin{equation}
        \Psi_N^*\Psi_N = \sum_{i=1}^v \left(\frac{s-\rho}{s+\rho}\right)^{i}\left(\frac{s+\rho}{s-\rho}\right)^i = v
    \end{equation}
\end{example}
